%%
%% ACM sigconf template for systems conferences (SOSP, ASPLOS, EuroSys).
%%
%% Submission mode: anonymous,review
%% Camera-ready: remove anonymous,review and fill in the \acm* rights block
%%               provided by ACM after acceptance.
%%

%% --- For submission (double-blind) ---
\documentclass[sigconf,anonymous,review]{acmart}

%% --- For camera-ready (uncomment and fill in after acceptance) ---
%\documentclass[sigconf]{acmart}
%\setcopyright{acmcopyright}
%\copyrightyear{2025}
%\acmYear{2025}
%\acmDOI{10.1145/XXXXXXX.XXXXXXX}
%\acmConference[EuroSys '25]{...}{...}{...}
%\acmBooktitle{...}
%\acmISBN{978-1-4503-XXXX-X/25/XX}

%% ============================================================
%% Packages
%% ============================================================
\usepackage{booktabs}    % Professional tables (\toprule, \midrule, \bottomrule)
\usepackage{graphicx}    % \includegraphics
\usepackage{subcaption}  % Subfigures
\usepackage{listings}    % Code listings
\usepackage{xspace}      % Smart spacing after macros
\usepackage{microtype}   % Better text justification
\usepackage{balance}     % Balance last page columns

%% ============================================================
%% Custom macros — define your system name once, use everywhere
%% ============================================================
\newcommand{\sysname}{SysName\xspace}   % Replace with your system's name
\newcommand{\eg}{e.g.,\xspace}
\newcommand{\ie}{i.e.,\xspace}
\newcommand{\etal}{\textit{et al.}\xspace}

%% ============================================================
%% Bibliography
%% ============================================================
\bibliographystyle{ACM-Reference-Format}
%% Use \cite{key} in text; put entries in references.bib

%%
%% ============================================================
\begin{document}

%% ============================================================
%% Title — anonymized for submission; restore authors for camera-ready
%% ============================================================
\title{\sysname: A [Noun Phrase Capturing Your Key Contribution]}

%% Submission: leave \author{} empty (anonymous mode handles it)
%% Camera-ready: fill in author names and affiliations
\author{Author One}
\affiliation{%
  \institution{Institution}
  \city{City}
  \country{Country}
}
\email{author1@example.com}

\author{Author Two}
\affiliation{%
  \institution{Institution}
  \city{City}
  \country{Country}
}
\email{author2@example.com}

%% ============================================================
%% Abstract
%% ============================================================
\begin{abstract}
%% 5-sentence formula for systems papers:
%% 1. The problem (concrete, with a number if possible)
%% 2. Why existing systems fail to solve it
%% 3. Key insight + high-level design approach
%% 4. Evaluation setup: workload(s), baseline(s), scale
%% 5. Headline result

[1. Concrete problem statement.] [2. Why existing systems fail.] [3. Key insight:
we observe that X, which enables our design to do Y.] [4. We implement \sysname
in N lines of C/Go/Rust and evaluate it on [workload] against [baseline] on a
[N-node cluster].] [5. \sysname achieves [Nx] improvement over [baseline] on
[benchmark], while reducing [overhead] by [Y\%].]
\end{abstract}

%% ============================================================
%% CCS Concepts and Keywords (required for camera-ready; optional for submission)
%% ============================================================
%% Generate at: https://dl.acm.org/ccs
%\begin{CCSXML}
%<ccs2012>
%   <concept>
%       <concept_id>10011007.10011006.10011008</concept_id>
%       <concept_desc>Software and its engineering~General programming languages</concept_desc>
%       <concept_significance>500</concept_significance>
%   </concept>
%</ccs2012>
%\end{CCSXML}
%\ccsdesc[500]{Software and its engineering~General programming languages}
%\keywords{keyword1, keyword2, keyword3}

\maketitle

%% ============================================================
\section{Introduction}
%% ============================================================

%% Opening scenario: a concrete, relatable situation where existing systems fail.
%% Include a number. "Consider a [system] serving [N] [operations]. When [event],
%% [consequence] — a [Nx] increase over normal [metric]."

[Opening scenario with concrete numbers showing existing systems fail.]

%% The fundamental challenge: why the problem is hard structurally.
[The challenge: why existing approaches fall short. What is the root cause?]

%% The key insight (explicit sentence).
\textbf{Key insight.} We observe that [one sentence insight that makes your
design possible]. This insight enables [high-level design description].

%% Brief design overview.
We design and implement \sysname, [one-sentence description of what it is and
what it does]. [Point to design section: §\ref{sec:design} describes the design.]

%% Contribution bullets (3–5, specific and falsifiable).
This paper makes the following contributions:
\begin{itemize}
  \item We identify [problem/observation] and show that [evidence] (\S\ref{sec:motivation}).
  \item We design \sysname, which achieves [goal] by [mechanism] (\S\ref{sec:design}).
  \item We implement \sysname in [N]k lines of [language] (\S\ref{sec:impl}).
  \item We show \sysname achieves [Nx] improvement over [baseline] on [benchmark]
        across [scale] (\S\ref{sec:eval}).
\end{itemize}

%% Evaluation summary.
We evaluate \sysname on [workloads] against [baselines] on a [N-node cluster].
\sysname achieves [headline result 1] and [headline result 2].

%% Paper organization.
\S\ref{sec:motivation} presents a motivating measurement study.
\S\ref{sec:design} describes the design of \sysname.
\S\ref{sec:impl} describes the implementation.
\S\ref{sec:eval} presents our evaluation.
\S\ref{sec:related} discusses related work.
\S\ref{sec:conclusion} concludes.


%% ============================================================
\section{Motivation}
\label{sec:motivation}
%% ============================================================

%% This section must:
%% 1. Show existing systems fail on a concrete workload (with measurement graphs)
%% 2. Explain WHY they fail structurally
%% 3. Connect the failure to your key insight

\subsection{The Problem in Practice}

[Describe the concrete scenario. Include Figure~\ref{fig:motivation} showing the
measured failure of existing systems.]

\begin{figure}[t]
  \centering
  \includegraphics[width=\columnwidth]{figs/motivation.pdf}
  \caption{[Existing system X] under [workload Y]: [metric] degrades from
           [A] to [B] when [condition], exceeding the [SLO] for [Z\%] of requests.}
  \label{fig:motivation}
\end{figure}

\subsection{Why Existing Systems Fall Short}

[Explain the structural reason existing systems fail—not just that they are slow,
but WHY their design cannot avoid this problem.]

\subsection{Key Insight}

[Restate the insight clearly. Connect it to the failure mode described above.]


%% ============================================================
\section{Design}
\label{sec:design}
%% ============================================================

%% State design goals and non-goals up front.
\sysname is designed to achieve the following goals:
\begin{enumerate}
  \item \textbf{Goal 1.} [Specific, measurable goal]
  \item \textbf{Goal 2.} [Specific, measurable goal]
  \item \textbf{Goal 3.} [Specific, measurable goal]
\end{enumerate}

\sysname does \emph{not} aim to [non-goal 1] or [non-goal 2].

\subsection{Overview}

[High-level architecture description with Figure~\ref{fig:arch}.]

\begin{figure}[t]
  \centering
  \includegraphics[width=\columnwidth]{figs/architecture.pdf}
  \caption{\sysname architecture. [Component A] handles [responsibility X];
           [Component B] manages [responsibility Y].
           Shaded components are new; unshaded are existing infrastructure.}
  \label{fig:arch}
\end{figure}

\subsection{Component A}
\label{sec:design:component-a}

[Detail Component A. Explain its interface, invariants, and key decisions.
Discuss important trade-offs made.]

\subsection{Component B}
\label{sec:design:component-b}

[Detail Component B.]

\subsection{Correctness}
\label{sec:design:correctness}

[For systems with strong correctness requirements, argue or sketch proof of
key safety/liveness properties. For engineering systems, discuss failure handling.]


%% ============================================================
\section{Implementation}
\label{sec:impl}
%% ============================================================

We implement \sysname in [language], comprising approximately [N] lines of code.
[System] runs on [OS, kernel version].

\paragraph{[Key implementation aspect 1.]}
[Describe it concisely. Mention specific files, modules, or kernel interfaces
if relevant.]

\paragraph{[Key implementation aspect 2.]}
[Describe it.]

\paragraph{Limitations of the current implementation.}
[Be honest about what's prototype-quality vs. production-ready.]


%% ============================================================
\section{Evaluation}
\label{sec:eval}
%% ============================================================

%% Every claim from §1 contribution bullets must have a corresponding experiment.

\subsection{Experimental Setup}

\paragraph{Testbed.}
[CPU model, cores, NUMA topology; Memory capacity and speed; Storage (NVMe/HDD,
model, capacity); Network (speed, topology); OS and kernel version.]

\paragraph{Baselines.}
We compare \sysname against:
\begin{itemize}
  \item \textbf{[Baseline 1]} [version]: [brief description and why it is the right comparison]
  \item \textbf{[Baseline 2]} [version]: [brief description]
\end{itemize}

\paragraph{Workloads.}
[Describe each workload used: benchmark name, configuration, parameters.]

\subsection{Main Results}
\label{sec:eval:main}

%% State the claim before showing the evidence.
\textbf{Claim:} \sysname achieves [Nx] higher throughput than [baseline] on [workload].

[Describe Figure~\ref{fig:main-result}. Tell the reader what to observe and why
it demonstrates the claim: "The blue line (\sysname) stays below [X~µs] even at
[Y~Mops/s], while [Baseline] exceeds [Z~µs] at [W~Mops/s] due to [reason]."]

\begin{figure}[t]
  \centering
  \includegraphics[width=\columnwidth]{figs/main-result.pdf}
  \caption{Throughput vs.\ p99 latency on [workload].
           \sysname achieves [N]$\times$ higher throughput at the [X~ms] p99 target.
           Error bars show one standard deviation over 5 runs.}
  \label{fig:main-result}
\end{figure}

\subsection{Scalability}
\label{sec:eval:scale}

[Show how the system scales with [cores / nodes / data size].]

\subsection{Overhead}
\label{sec:eval:overhead}

[Measure \sysname's overhead on the unmodified workload / common case.]

\subsection{Microbenchmarks}
\label{sec:eval:micro}

[Microbenchmarks isolating specific components. Tie each back to a claim.]

\subsection{Summary}

[1-2 sentence takeaway of what the evaluation demonstrates.]


%% ============================================================
\section{Related Work}
\label{sec:related}
%% ============================================================

%% Organize THEMATICALLY, not paper-by-paper.
%% Good: "A line of work [refs] handles X by Y; we differ because Z."
%% Bad: "Smith et al. did A. Jones et al. did B."

\paragraph{[Theme 1: e.g., KV stores].}
[Group of papers doing X. How they differ from \sysname and why \sysname's
approach is needed.]

\paragraph{[Theme 2].}
[...]

\paragraph{[Theme 3].}
[...]


%% ============================================================
\section{Conclusion}
\label{sec:conclusion}
%% ============================================================

We present \sysname, [one sentence summary of what it is and what it achieves].
Our key insight is that [restate insight]. \sysname achieves [headline result]
over [baseline] on [workload].

[Optional: 1-2 sentences on future work or broader implications.]


%% ============================================================
%% Acknowledgments (remove or anonymize for submission)
%% ============================================================
%\begin{acks}
%We thank [...]
%\end{acks}


%% ============================================================
%% References
%% ============================================================
\bibliography{references}

\end{document}
